\documentclass[11pt]{article}

\usepackage[T1]{fontenc}
\usepackage{amsmath,amssymb,amsthm,mathtools}
\usepackage[margin=1in]{geometry}
\usepackage{microtype}

\newtheorem{theorem}{Theorem}
\newtheorem{lemma}[theorem]{Lemma}
\newtheorem{proposition}[theorem]{Proposition}
\newtheorem{claim}[theorem]{Claim}

\newcommand{\PP}{\mathbb P}
\newcommand{\EE}{\mathbb E}
\newcommand{\cF}{\mathcal F}
\newcommand{\cA}{\mathcal A}
\newcommand{\cB}{\mathcal B}
\newcommand{\1}{\mathbf 1}

\title{Canonical Optimal Wallets Have at Most One Incomplete Layer}
\author{}
\date{}

\begin{document}

\maketitle

\section{Model and statement}

Let $[n]=\{1,\dots,n\}$.  Independently for each key, let its state be
\[
\begin{array}{c|c|c|c}
\text{state}&\text{owner has the key}&\text{attacker has the key}
&\text{probability}\\
\hline
\text{safe}&1&0&s\\
\text{stolen}&0&1&t\\
\text{leaked}&1&1&e\\
\text{lost}&0&0&\ell
\end{array}
\]
where
\[
s+t+e+\ell=1,
\qquad
s>\frac12.
\]
Write $O,D\subseteq[n]$ for the sets of keys held by the owner and the
attacker, respectively.  A \emph{wallet} is an upward-closed family
$\cF\subseteq 2^{[n]}$.  Its success probability is
\[
\Phi_{s,t,e,\ell}(\cF)
   =\PP(O\in\cF,\ D\notin\cF).
\]
When the parameters are fixed, we abbreviate this to $\Phi(\cF)$.

For $0\leq k\leq n$, let
\[
\cF_k=\cF\cap\binom{[n]}k.
\]
The $k$th layer is called \emph{incomplete} if
\[
0<|\cF_k|<\binom nk.
\]

\begin{theorem}\label{thm:main}
For every wallet $\cF\subseteq2^{[n]}$, there is a wallet
$\cF^\star\subseteq2^{[n]}$ such that
\[
\Phi(\cF^\star)\geq\Phi(\cF)
\]
and $\cF^\star$ has at most one incomplete layer.  Consequently, some
optimal wallet has at most one incomplete layer.
\end{theorem}

The proof uses three batch edits at the two extreme incomplete layers.
The main ingredient is a rank-kernel inequality that remains valid in
the presence of leaked keys.

\section{The three batch edits}

Suppose that $\cF$ has at least two incomplete layers, and let $r<q$ be
the lowest and highest incomplete ranks.  Upward closure implies that
every layer below $r$ is empty and every layer above $q$ is full.  Put
\[
\cA=\cF_r,
\qquad
\cB=\binom{[n]}q\setminus\cF_q,
\]
and define
\[
\alpha=\frac{|\cA|}{\binom nr},
\qquad
\beta=\frac{|\cB|}{\binom nq}.
\]
Both $\alpha$ and $\beta$ are strictly between zero and one.  Consider
the three wallets
\[
\cF_R=\cF\setminus\cA,
\qquad
\cF_Q=\cF\cup\cB,
\qquad
\cF_{RQ}=(\cF\setminus\cA)\cup\cB.
\]
They are upward closed.  Indeed, removing the lowest nonempty layer
cannot violate upward closure, and filling the highest non-full layer
cannot violate it.

Let
\[
I=|O|,
\qquad
J=|D|,
\]
and, for integers $k,m$, define
\begin{equation}\label{eq:def-g}
g_k(m)=\PP(I\geq m,J=k)-\PP(I=k,J<m).
\end{equation}

\begin{lemma}[Batch-edit bounds]\label{lem:batch-bounds}
The changes caused by the three edits satisfy
\begin{align}
\Phi(\cF_R)-\Phi(\cF)
   &\geq \alpha g_r(q+1),\label{eq:remove-bound}\\
\Phi(\cF_Q)-\Phi(\cF)
   &\geq-\beta g_q(r),\label{eq:fill-bound}\\
\Phi(\cF_{RQ})-\Phi(\cF)
   &\geq L_1:=\alpha g_r(q)-\beta g_q(r),\label{eq:joint-bound-1}\\
\Phi(\cF_{RQ})-\Phi(\cF)
   &\geq L_2:=\alpha g_r(q+1)-\beta g_q(r+1).
   \label{eq:joint-bound-2}
\end{align}
\end{lemma}

\begin{proof}
Conditioned on $I=k$, the owner set $O$ is uniform on
$\binom{[n]}k$; conditioned on $J=k$, the attacker set $D$ is uniform
on $\binom{[n]}k$.  The same statements remain true after imposing any
condition involving only $(I,J)$.  This follows directly from the
exchangeability of the $n$ coordinates.

For the remove edit, a success is certainly gained whenever
$D\in\cA$ and $I\geq q+1$, because all layers above $q$ are full.  A
success can be lost only when $O\in\cA$ and $J\leq q$, because an
attacker with at least $q+1$ keys was already authorized.  Uniformity
within a rank therefore gives
\[
\Phi(\cF_R)-\Phi(\cF)
\geq
\alpha\bigl[\PP(I\geq q+1,J=r)-\PP(I=r,J<q+1)\bigr],
\]
which is \eqref{eq:remove-bound}.

For the fill edit, a success is certainly gained whenever $O\in\cB$
and $J<r$, because every layer below $r$ is empty.  A success can be
lost only when $D\in\cB$ and $I\geq r$, because every authorized owner
has at least $r$ keys.  Thus
\[
\Phi(\cF_Q)-\Phi(\cF)
\geq
\beta\bigl[\PP(I=q,J<r)-\PP(I\geq r,J=q)\bigr],
\]
which is \eqref{eq:fill-bound}.

There are two ways to perform the joint edit.  If layer $q$ is filled
first, the fill step contributes at least $-\beta g_q(r)$.  In the
resulting wallet, layer $q$ is full, so the subsequent removal of
$\cA$ contributes at least $\alpha g_r(q)$.  This proves
\eqref{eq:joint-bound-1}.  If $\cA$ is removed first, the removal step
contributes at least $\alpha g_r(q+1)$.  In the resulting wallet every
layer through rank $r$ is empty, so subsequently filling $\cB$
contributes at least $-\beta g_q(r+1)$.  This proves
\eqref{eq:joint-bound-2}.
\end{proof}

The elementary discrete derivative
\begin{equation}\label{eq:g-monotone}
g_k(m+1)-g_k(m)
=-\PP(I=m,J=k)-\PP(I=k,J=m)
\leq0
\end{equation}
will be used repeatedly.

\section{The low-sum kernel certificate}

\begin{proposition}[Kernel certificate]\label{prop:certificate}
If
\[
1\leq r<q\leq n-1,
\qquad
r+q\leq n,
\]
then
\begin{equation}\label{eq:certificate}
\boxed{\binom nq\,g_r(q+1)
\geq
\binom nr\,g_q(r+1).}
\end{equation}
\end{proposition}

We first show that this certificate implies the desired batch
improvement when $r+q\leq n$.

\begin{lemma}[Low-sum batch improvement]\label{lem:low-sum}
If $r+q\leq n$, at least one of $\cF_R,\cF_Q,\cF_{RQ}$ has success
probability at least $\Phi(\cF)$.
\end{lemma}

\begin{proof}
If $g_r(q+1)\geq0$, \eqref{eq:remove-bound} shows that the remove edit
is non-worsening.  If $g_q(r)\leq0$, \eqref{eq:fill-bound} shows that
the fill edit is non-worsening.  It remains to consider
\begin{equation}\label{eq:crossing-regime}
g_r(q+1)<0<g_q(r).
\end{equation}
Set
\[
w=\PP(I=q,J=r)+\PP(I=r,J=q)
\]
and
\[
\rho=\frac{\binom nq}{\binom nr}.
\]
Since $r<q\leq n-r$, binomial unimodality and symmetry give
$\rho\geq1$.  By \eqref{eq:g-monotone},
\begin{equation}\label{eq:two-jumps}
g_r(q+1)=g_r(q)-w,
\qquad
g_q(r+1)=g_q(r)-w.
\end{equation}
The certificate \eqref{eq:certificate} says
\[
\rho\bigl(g_r(q)-w\bigr)\geq g_q(r)-w.
\]
Rearranging,
\[
g_r(q)-g_q(r)
\geq
(\rho-1)\bigl(w-g_r(q)\bigr).
\]
The first inequality in \eqref{eq:crossing-regime} and
\eqref{eq:two-jumps} imply $w>g_r(q)$, so
\begin{equation}\label{eq:cross-order-1}
g_r(q)\geq g_q(r)>0.
\end{equation}
Using \eqref{eq:two-jumps} again gives
\begin{equation}\label{eq:cross-order-2}
g_q(r+1)\leq g_r(q+1)<0.
\end{equation}

If $L_1\geq0$, \eqref{eq:joint-bound-1} proves that the joint edit is
non-worsening.  Otherwise
\[
\alpha g_r(q)<\beta g_q(r).
\]
By \eqref{eq:cross-order-1}, this implies $\alpha<\beta$.  By
\eqref{eq:cross-order-2},
\[
-g_q(r+1)\geq-g_r(q+1)>0.
\]
Consequently
\[
\begin{aligned}
L_2
&=\beta[-g_q(r+1)]-\alpha[-g_r(q+1)]\\
&\geq(\beta-\alpha)[-g_r(q+1)]>0.
\end{aligned}
\]
Now \eqref{eq:joint-bound-2} proves that the joint edit is improving.
\end{proof}

It remains to prove Proposition~\ref{prop:certificate}.  We do this
coefficientwise.

\section{Coefficient reduction}

Put
\begin{equation}\label{eq:eta-substitution}
\eta=s-(t+e+\ell)=2s-1\geq0,
\qquad
s=t+e+\ell+\eta.
\end{equation}
After this substitution, the difference between the left- and
right-hand sides of \eqref{eq:certificate} is a homogeneous polynomial
in $t,e,\ell,\eta$.  We shall prove that every coefficient is
nonnegative.

Fix a monomial
\[
t^a e^b\ell^c\eta^d,
\qquad
a+b+c+d=n,
\]
and fix a labelled partition
\[
[n]=T\sqcup E\sqcup L\sqcup H,
\qquad
|T|=a,\quad |E|=b,\quad |L|=c,\quad |H|=d.
\]
In the expansion \eqref{eq:eta-substitution}, this partition has the
following interpretation:
\[
\begin{array}{c|c}
\text{block}&\text{allowed states}\\
\hline
T&\text{safe or stolen}\\
E&\text{safe or leaked}\\
L&\text{safe or lost}\\
H&\text{safe only}.
\end{array}
\]
Call a choice of one of the allowed states at every coordinate a
\emph{toggle assignment}.  If $x,y,z$ denote the numbers chosen as
stolen, leaked, and lost in the first three blocks, respectively, then
\begin{equation}\label{eq:toggle-ranks}
J=x+y,
\qquad
n-I=x+z,
\end{equation}
and the multiplicity of these choices is
\[
\binom ax\binom by\binom cz.
\]

For this fixed labelled partition define
\begin{equation}\label{eq:def-G}
G_k(m)
=\#\{J=k,\ n-I\leq n-m\}
-\#\{n-I=n-k,\ J<m\},
\end{equation}
where both counts range over all toggle assignments.  The coefficient
of $t^ae^b\ell^c\eta^d$ in the polynomial from
\eqref{eq:certificate} is the common positive multinomial factor
\[
\frac{n!}{a!\,b!\,c!\,d!}
\]
times
\begin{equation}\label{eq:coefficient-target}
\binom nqG_r(q+1)-\binom nrG_q(r+1).
\end{equation}
Thus Proposition~\ref{prop:certificate} follows once
\eqref{eq:coefficient-target} is shown to be nonnegative for every
labelled partition.

\section{Averaging over nested sets}

For $A\subset B\subseteq[n]$, with $|A|=r$ and $|B|=q$, define
\begin{equation}\label{eq:Lambda}
\begin{aligned}
\Lambda(A,B)={}&
\1_{\{D=A,\ |O|\geq q+1\}}
+\1_{\{O=B,\ |D|\leq r\}}\\
&-\1_{\{O=A,\ |D|\leq q\}}
-\1_{\{D=B,\ |O|\geq r+1\}}.
\end{aligned}
\end{equation}
Summing over all toggle assignments and all nested pairs gives
\begin{align}
&\sum_{\substack{A\subset B\\|A|=r,\ |B|=q}}
\ \sum_{\text{toggle assignments}}\Lambda(A,B)
\notag\\
&\qquad
=\binom{n-r}{q-r}G_r(q+1)
-\binom qrG_q(r+1)\notag\\
&\qquad
=\frac{\binom qr}{\binom nr}
\left[
\binom nqG_r(q+1)-\binom nrG_q(r+1)
\right].
\label{eq:nested-average}
\end{align}
Indeed, a fixed $r$-set occurs as the lower member of
$\binom{n-r}{q-r}$ nested pairs, while a fixed $q$-set contains
$\binom qr$ lower $r$-sets; the identity in the last line follows by
double-counting pairs $(A,B)$ with $A\subset B$.

We prove that the left-hand side of \eqref{eq:nested-average} is
nonnegative.  Fix
\[
M=B\setminus A,
\qquad
|M|=q-r,
\]
and sum over all $r$-sets $A$ disjoint from $M$.  There are four
exhaustive types.

\begin{enumerate}
\item If $M$ contains an $H$-coordinate, both negative events in
\eqref{eq:Lambda} are impossible: an $H$-coordinate belongs to $O$ and
does not belong to $D$ in every toggle assignment.

\item If $M$ contains both an $E$-coordinate and an $L$-coordinate,
then $O=A$ is impossible because every $E$-coordinate belongs to $O$,
whereas $D=B$ is impossible because no $L$-coordinate belongs to $D$.

\item If every coordinate of $M$ lies in $T$, switching all
coordinates of $M$ from stolen to safe gives injections
\[
\{O=A,\ |D|\leq q\}
\longrightarrow
\{O=B,\ |D|\leq r\}
\]
and
\[
\{D=B,\ |O|\geq r+1\}
\longrightarrow
\{D=A,\ |O|\geq q+1\}.
\]
The inverse on either image switches the coordinates of $M$ back from
safe to stolen, so both maps are injective.

\item The only remaining types are
\[
M\subseteq T\cup L,\quad M\cap L\neq\varnothing,
\]
and
\[
M\subseteq T\cup E,\quad M\cap E\neq\varnothing.
\]
\end{enumerate}

The next two sections handle these final two types.

\section{The $L$-case}

Assume
\[
M\subseteq T\cup L,
\qquad
l:=|M\cap L|\geq1.
\]
Outside $M$, let the numbers of $T,E,L,H$ coordinates be
\[
T_0,\quad b,\quad L_0,\quad d,
\]
respectively.  Put
\[
\delta=n-r-q\geq0,
\qquad
k_0=r-b-d.
\]
Since $|M|=q-r$, elementary counting gives
\begin{equation}\label{eq:L-identities}
T_0=2r-b-L_0-d+\delta,
\qquad
2k_0=T_0-b+L_0-d-\delta.
\end{equation}

We first record the event count that will be used.  In the negative
event $O=A$, every outside-$M$ coordinate of $E\cup H$ is forced into
$A$.  If $A$ contains $j$ of the outside $L$-coordinates, it must
contain $k_0-j$ outside $T$-coordinates.  Pairing this event with the
positive event $O=B$ leaves precisely the configurations in which the
number of leaked outside $E$-coordinates is
\[
L_0-\delta-j+s,
\qquad
1\leq s\leq l.
\]

For the other positive event, $D=A$, reindex by
\[
|A\cap E|=L_0-\delta-j.
\]
The required $T$-choice has size complementary to $k_0-j$, and
therefore has the same binomial count $\binom{T_0}{k_0-j}$.  If $W$ is
the set of lost coordinates in the combined $L$-block of size
$L_0+l$, the condition $|O|\geq q+1$ is exactly
\[
|W|\leq L_0-j-1.
\]
Consequently, after summing \eqref{eq:Lambda} over all $A$ disjoint
from this fixed $M$, the desired nonnegativity is equivalent to
\begin{equation}\label{eq:L-sum}
\sum_{j\in\mathbb Z}\binom{T_0}{k_0-j}H_j\geq0,
\end{equation}
where
\begin{equation}\label{eq:def-Hj}
\begin{aligned}
H_j={}&
\binom b{L_0-\delta-j}
\sum_{w=0}^{L_0-j-1}\binom{L_0+l}{w}\\
&-\binom{L_0}j
\sum_{s=1}^{l}\binom b{L_0-\delta-j+s}.
\end{aligned}
\end{equation}
As usual, an out-of-range binomial coefficient and an empty sum are
zero.  In particular, $j$ ranges over all integers.

We prove \eqref{eq:L-sum} by an injection from the objects counted by
the negative part to those counted by the positive part.  A negative
object indexed by $(j,s)$ consists of
\[
\begin{aligned}
K&\subseteq[T_0],& |K|&=k_0-j=:k,\\
J&\subseteq[L_0],& |J|&=j,\\
A_E&\subseteq[b],& |A_E|&=L_0-\delta-j+s,\\
&&1&\leq s\leq l.
\end{aligned}
\]
Set
\[
j'=j-s.
\]
Keep $A_E$, and let
\begin{equation}\label{eq:L-W-map}
W=([L_0]\setminus J)\cup P_{s-1},
\end{equation}
where $P_{s-1}$ is the canonical prefix of length $s-1$ in the fixed,
ordered $l$-element set $M\cap L$.  Then
\[
|W|=L_0-j+s-1=L_0-j'-1,
\]
so $W$ is counted by the cumulative sum in the positive term of
\eqref{eq:def-Hj} at index $j'$.

It remains to raise $K$ from size $k$ to size $k+s$.  Nonvanishing of
the negative object implies
\[
L_0-\delta-j+s\leq b.
\]
Using \eqref{eq:L-identities}, we obtain
\begin{align}
T_0-2k-s
&=b-L_0+d+\delta+2j-s\notag\\
&=(b-L_0+\delta+j-s)+(d+j)\geq0.
\label{eq:L-raising-condition}
\end{align}
Hence $k+s\leq T_0-k$.  Fix a symmetric-chain decomposition of
$2^{[T_0]}$.  Moving upward by $s$ positions in the chain containing
$K$ defines an injective map
\[
\binom{[T_0]}k\longrightarrow\binom{[T_0]}{k+s};
\]
condition \eqref{eq:L-raising-condition} guarantees that the target
position exists.  Let $K'$ be the image of $K$.

The resulting triple $(K',A_E,W)$ is a positive object at index $j'$.
It recovers $s$ from the length of the canonical prefix
$W\cap(M\cap L)$, then recovers $j$, $J$, and $A_E$, and finally
recovers $K$ from the injective raising map.  Thus the map is
injective, and images belonging to different $(j,s)$ are disjoint.
This proves \eqref{eq:L-sum}.

\section{The $E$-case}

Assume
\[
M\subseteq T\cup E,
\qquad
e_M:=|M\cap E|\geq1.
\]
Outside $M$, let the numbers of $T,E,L,H$ coordinates be
\[
T_0,\quad E_0,\quad c,\quad d,
\]
respectively, and put
\[
\delta=n-r-q,
\qquad
k_0=r-E_0-d.
\]

We again first record the event count.  In the positive event $O=B$,
suppose $A$ contains $j$ outside $L$-coordinates.  It must contain
every outside $E$- and $H$-coordinate, leaving $k_0-j$ choices from
the outside $T$-block.  The condition $|D|\leq r$ permits at most
\[
c-\delta-j
\]
leaked coordinates in the combined $E$-block of size $E_0+e_M$.

For the negative event $D=B$, first pair all configurations covered by
the positive event $D=A$ by forcing the coordinates of $M$ safe.
Reindex the remaining configurations by
\[
|A\cap E|=c-\delta-j.
\]
The unmatched configurations have an outside-$L$ safe set of size
\[
j+1-s,
\qquad
1\leq s\leq e_M.
\]
It follows that the fixed-$M$ sum is
\begin{equation}\label{eq:E-sum}
\begin{aligned}
\sum_{j\in\mathbb Z}\binom{T_0}{k_0-j}
\Bigg[
&\binom cj\sum_{w=0}^{c-\delta-j}\binom{E_0+e_M}{w}\\
&-\binom{E_0}{c-\delta-j}
\sum_{s=1}^{e_M}\binom c{j+1-s}
\Bigg].
\end{aligned}
\end{equation}

A negative object indexed by $(j,s)$ consists of
\[
\begin{aligned}
K&\subseteq[T_0],& |K|&=k_0-j=:k,\\
A_E&\subseteq[E_0],& |A_E|&=c-\delta-j,\\
J&\subseteq[c],& |J|&=j+1-s,\\
&&1&\leq s\leq e_M.
\end{aligned}
\]
Set
\[
j'=j+1-s.
\]
Keep $J$, extend $A_E$ by the canonical prefix of length $s-1$ in
the ordered set $M\cap E$, and raise $K$ by $s-1$ ranks in the fixed
symmetric-chain decomposition of $2^{[T_0]}$.

The required raising map exists.  Indeed, nonvanishing of $A_E$ and
$J$ gives
\[
E_0-c+\delta+j\geq0,
\qquad
j-s+1\geq0.
\]
The analogue of \eqref{eq:L-identities} therefore yields
\begin{align}
T_0-2k-(s-1)
&=(E_0-c+\delta+j)+(d+j-s+1)\geq0.
\label{eq:E-raising-condition}
\end{align}
Thus $k+s-1\leq T_0-k$, as required.  The enlarged $E$-set has size
\[
c-\delta-j+s-1=c-\delta-j',
\]
which is the top allowed size in the positive cumulative term of
\eqref{eq:E-sum} at index $j'$.

The canonical prefix records $s$.  Hence the image recovers $j$,
$A_E$, and $J$, and injectivity of the symmetric-chain raising map
recovers $K$.  Images from different $(j,s)$ are disjoint.  Therefore
the negative objects in \eqref{eq:E-sum} inject into the positive
objects, so \eqref{eq:E-sum} is nonnegative.

Every fixed-$M$ contribution in \eqref{eq:nested-average} is now
nonnegative.  It follows that \eqref{eq:coefficient-target} is
nonnegative for every labelled partition.  Every coefficient of the
polynomial obtained after \eqref{eq:eta-substitution} is therefore
nonnegative.  Since $t,e,\ell,\eta\geq0$, the polynomial itself is
nonnegative.  This proves Proposition~\ref{prop:certificate}.

\section{The high-sum case by duality}

\begin{lemma}[High-sum batch improvement]\label{lem:high-sum}
If $r+q>n$, at least one of $\cF_R,\cF_Q,\cF_{RQ}$ has success
probability at least $\Phi(\cF)$.
\end{lemma}

\begin{proof}
Define the dual wallet
\[
\cF^\ast
=\{X\subseteq[n]:[n]\setminus X\notin\cF\}.
\]
It is upward closed.  Define transformed owner and attacker sets by
\[
O^\ast=[n]\setminus D,
\qquad
D^\ast=[n]\setminus O.
\]
Coordinatewise, this transformation preserves safe and stolen states
and exchanges leaked and lost states.  Moreover,
\[
O^\ast\in\cF^\ast,\ D^\ast\notin\cF^\ast
\quad\Longleftrightarrow\quad
O\in\cF,\ D\notin\cF.
\]
Consequently
\begin{equation}\label{eq:duality-phi}
\Phi_{s,t,e,\ell}(\cF)
=\Phi_{s,t,\ell,e}(\cF^\ast).
\end{equation}

The lowest and highest incomplete ranks of $\cF^\ast$ are
\[
r^\ast=n-q,
\qquad
q^\ast=n-r.
\]
Since $r+q>n$,
\[
r^\ast+q^\ast=2n-r-q<n.
\]
Lemma~\ref{lem:low-sum}, applied to $\cF^\ast$ with the leaked and lost
probabilities interchanged, gives a non-worsening remove, fill, or
joint edit of $\cF^\ast$.  Under duality, removal at the lowest
incomplete layer becomes filling at the highest incomplete layer of
$\cF$, filling becomes removal, and the joint edit remains joint.
Equation \eqref{eq:duality-phi} therefore gives a non-worsening batch
edit of $\cF$.
\end{proof}

\section{Completion of the proof}

\begin{proof}[Proof of Theorem~\ref{thm:main}]
If $\cF$ already has at most one incomplete layer, take
$\cF^\star=\cF$.  Otherwise, let $r<q$ be its extreme incomplete
ranks.  Lemma~\ref{lem:low-sum} or Lemma~\ref{lem:high-sum}, according
as $r+q\leq n$ or $r+q>n$, supplies one of the three batch edits whose
success probability is at least $\Phi(\cF)$.

The remove edit makes layer $r$ empty, the fill edit makes layer $q$
full, and the joint edit does both.  No other layer changes.  Hence
each batch edit eliminates at least one incomplete layer and creates
none.  After finitely many steps, the process terminates at a wallet
$\cF^\star$ with at most one incomplete layer and
\[
\Phi(\cF^\star)\geq\Phi(\cF).
\]

Finally, begin with an optimal wallet.  Every non-worsening edit must
then preserve the optimal value.  The terminal wallet is therefore
optimal and has at most one incomplete layer.
\end{proof}

\end{document}
